% 第13节：III型因子与Split性质
\section{III型因子与Split性质}
\label{sec:type_III}

\subsection{引言}

量子场论中的局域代数通常是III$_1$型因子。本节讨论CTUFT的期望结构。

\subsection{类型分类}

\begin{definition}[因子类型]
具有平凡中心（$\mathcal{M} \cap \mathcal{M}' = \mathbb{C}\mathbb{I}$）的von Neumann代数$\mathcal{M}$是：
\begin{itemize}
    \item \textbf{I型}：具有极小投影（矩阵代数）
    \item \textbf{II型}：具有有限迹但没有极小投影
    \item \textbf{III型}：不存在迹
\end{itemize}
\end{definition}

\subsection{CTUFT局域代数}

\begin{conjecture}[III$_1$型性质]
对于任何开有界区域$\mathcal{O} \subset \mathbb{R}^4$，CTUFT中的局域代数$\mathcal{A}(\mathcal{O})$是III$_1$型因子。
\end{conjecture}

\textbf{证据：}
\begin{enumerate}
    \item 具有质量间隙的一般QFT具有III型局域代数
    \item 挠率场具有连续谱（没有极小投影）
    \item 模理论显示非平凡的$\Delta$谱
\end{enumerate}

\subsection{Split性质}

\begin{conjecture}[Split性质]
对于$\mathcal{O}_1 \subset \mathcal{O}_2$且有限分离，存在I型因子$\mathcal{N}$：
\begin{equation}
    \mathcal{A}(\mathcal{O}_1) \subset \mathcal{N} \subset \mathcal{A}(\mathcal{O}_2)
\end{equation}
\end{conjecture}

这一性质确保纠缠具有足够的局域自由度。

\subsection{当前状态}

III型结构和split性质对自由场已确立。对于相互作用的CTUFT：
\begin{itemize}
    \item 自由极限：已证明为III$_1$型
    \item 微扰：证据支持III型
    \item 非微扰：开放问题
\end{itemize}
